\documentclass[11pt]{article}
\usepackage[letterpaper,margin=1in]{geometry}
\usepackage{fontspec}
\usepackage{amsmath,amssymb}
\usepackage{graphicx,booktabs,xcolor,microtype}
\usepackage[numbers,sort&compress]{natbib}
\usepackage[font=small,labelfont=bf]{caption}
\usepackage[hidelinks,pdfusetitle]{hyperref}
\graphicspath{{figures/}}
\newcommand{\ShuffleSnap}{ShuffleSnap}
\newcommand{\panel}[2]{\par\begin{center}\begin{minipage}{\linewidth}\centering
  \includegraphics[width=\linewidth]{#1}\captionof{figure}{#2}\end{minipage}\end{center}}

\title{\ShuffleSnap: strided windowed exact reassignment\\for point-to-grid layout}
\author{Kyle McDonald}
\date{October 6, 2026}

\begin{document}
\maketitle

\begin{abstract}
Giving each point of a 2D embedding its own grid cell with the least total squared
displacement is a linear assignment problem too large for dense solvers at $10^5$ points.
\ShuffleSnap, inspired by Shellsort, solves small exact assignment problems in windows whose
cells lie a stride apart, and shrinks the stride to one. Against certified optima it comes
within hundredths of a percent on clustered layouts, including real UMAP embeddings, in
seconds. Its errors are loops of one-cell shifts that no window can see, which make uniform
points, concave masks and surplus cells hard.
\end{abstract}

\section{From points to cells}

An image collection is often shown by embedding the items in the plane, for example with
t-SNE or UMAP~\citep{vandermaaten2008tsne,mcinnes2018umap}, and placing each thumbnail in
its own grid cell. Rounding every point to its nearest cell causes collisions and holes, so
the cells must be assigned jointly: we seek the assignment $\sigma$ of points $p_i$ to
distinct cells $c_{\sigma(i)}$ that minimizes $\sum_i \lVert p_i - c_{\sigma(i)}\rVert^2$.
This linear assignment problem, also a discrete optimal-transport
problem~\citep{peyre2019computational}, is solved exactly by dense
solvers~\citep{kuhn1955hungarian,jonker1987shortest} that need the full cost matrix: 80\,GB
at $10^5$ points. Practical tools instead split the cloud recursively, like
RasterFairy~\citep{klingemann2015rasterfairy}, or merge randomized partitions, like
BSP-OT~\citep{genest2025bspot}.

\ShuffleSnap\ borrows from Shellsort~\citep{shell1959highspeed}, which sorts interleaved
subsequences at a shrinking gap. It solves only small assignment problems, but exactly,
with their cells spread across the grid first and brought together later
(Figure~\ref{fig:schematic}). We measure its gaps against a solver that proves optimality
on the same instances, and report where it fails as well as where it works.

\panel{fig_schematic}{(a)~A window collects the points that occupy its cells and
reassigns them among the same cells by an exact linear assignment. (b)~At stride $s$ a
window takes every $s$-th row and column: it keeps 64 cells but spans $7s+1$ cells per
axis, so points can travel far. Strides shrink from left to right. (c)~Each subgrid is
tiled four times, shifted by half a window, so that points can cross every window
border.\label{fig:schematic}}

\clearpage

\section{Coarse to fine}

\ShuffleSnap\ starts from a random assignment. At stride $s$ we keep every $s$-th row
and every $s$-th column of the grid, which splits it into $s^2$ interleaved subgrids, one
for each choice of first row and first column. Each subgrid is cut into windows of
$8\times8$ of its cells. In every window, the current occupants are reassigned among the
window's cells by solving a linear assignment problem of at most $64\times64$ exactly,
with costs taken from the original point coordinates. The new assignment replaces the old
one only if it lowers the window's cost by more than a relative $10^{-12}$, so the layout
stays valid, its cost never increases, and ties cannot cause cycles.

A round tiles every subgrid four times, shifted by $(0,0)$, $(0,4)$, $(4,0)$ and $(4,4)$.
Windows that cross an edge are clipped, not dropped. The windows of one tiling are disjoint, so they are solved in parallel. On an
axis of length $L$ the first stride is the smallest power of two with $8s\ge L$, which
makes the first windows span the grid. After each round the stride is divided by $\sqrt2$
 and rounded; at stride 1, rounds
repeat until one moves no point. Each level costs about the same, linear in the number of
cells, and the no-move stop is a fixed point of the windows, not a proof of optimality.
We chose the ratio $\sqrt2$ experimentally: no ratio between 1.1 and 3 was clearly better.
Larger windows ended closer to the optimum in every test, but took longer to reach any
cost that a smaller window also reached. Besides the default
$8\times8$ we offer $5\times5$ (\emph{fast}) and $12\times12$ (\emph{quality}).

\panel{fig_progress}{Progress on 10{,}000 points of a UMAP embedding of MNIST
digits~\citep{lecun1998gradient} (every second stride shown). Every point keeps the color
of its input position, and the dark bands are graph-paper lines drawn on the input: both
travel with the points. Numbers are the mean squared displacement in cells$^2$.\label{fig:progress}}

\paragraph{Targets and implementation.} A target can be any set of cells: a rectangle, a
disk, the shape of a letter. Each axis has its own stride, masked cells never enter a
window, and empty cells do, so with more cells than points the solver also chooses which
cells stay empty. The kernel is C++; a 70-line NumPy version reproduces it exactly.

\paragraph{Certified optima.} Measuring gaps needs the true optimum, but at $10^5$ points
a dense solver does not fit in memory. We solve a min-cost flow with
OR-Tools~\citep{ortools} on sparse candidate edges: each point's 12 nearest cells and the
neighbourhood of its cell in a valid starting layout (\ShuffleSnap's). As this may miss an
edge, we check it against the full, dense problem, whose optimal cost any cell potentials
$v$ bound from below:
\[
  \mathrm{LB}(v)=\sum_i \min_j\bigl(\lVert p_i-c_j\rVert^2-v_j\bigr)+\sum_j v_j-(M-N)\max_j v_j ,
\]
where $M\ge N$ is the number of cells. We derive $v$ from the sparse solution, evaluate
the bound over \emph{all} cells, add every edge it shows to be missing, and repeat until
the bound meets the cost. Then no assignment of the dense problem is cheaper: the result
is its exact optimum; the start affects only the time. Where it fits, SciPy's dense solver found the same optimum on 600 instances of up to 2{,}500 points and on all ten
distributions at $10^4$ points (relative difference below $10^{-9}$). It certified every benchmark instance (up to $2\times10^5$ points); all gaps below use it.

\clearpage

\section{Evidence and limits}

\panel{fig_summary}{Excess over the certified optimum against time on one core.
Points are means over three instances (for MNIST, three runs on the one embedding); bars
span their range. The three \ShuffleSnap\ presets use $5\times5$, $8\times8$ (default) and
$12\times12$ windows from left to right. The vertical line is the time of the certified exact
solver.\label{fig:summary}}

We ran ten synthetic distributions with $10^5$ points and the complete UMAP embeddings
of MNIST and Fashion-MNIST~\citep{xiao2017fashion} (70{,}000 points each) on one
workstation. Every method received identical coordinates and cells, and all times include
each method's own preparation. Figure~\ref{fig:summary} shows three of them and MNIST.

On clustered data, the kind of layout that embeddings produce, \ShuffleSnap\ is close to
optimal. It is within 0.06\% of the optimum on nine of the ten distributions and takes
a median of 2.3\,s on one core, or 0.19\,s on 24 threads. On the two real embeddings it is
within 0.002\% of the optimum, and its layouts are visually indistinguishable from the
optimal ones (Figure~\ref{fig:visual}). At equal time it beats BSP-OT on most inputs: the fast preset (0.77\,s) is
closer to the optimum than BSP-OT with 16 plans (0.71\,s) on eight of ten distributions,
and the default (2.3\,s) than 64 plans (2.9\,s) on nine; on uniform points they are at 2.6\%
and 0.84\%, against 45\% and 32\%. RasterFairy is 17--425\% above the optimum on seven. Displacement is not the
whole story, though: on the real embeddings BSP-OT stays within a few percent of the
optimal displacement yet halves the fraction of grid neighbours that were input
neighbours, because it scrambles local order.

\ShuffleSnap\ is a local search, and its failures follow from that. Where it misses the
optimum, the difference forms a few long closed loops in which every point is shifted by
about one cell. Such a loop is optimal inside every window that sees only part of it, and
a strided window cannot shift a point by one cell. We built inputs on which the method
never moves while the optimum costs zero, so its relative excess is unbounded. On uniform
points the loops grow with $N$: the excess rises from 0.1\% at $10^3$ points to 0.84\% at
$10^5$ and 1.3\% at $10^6$. Concave masks and grids with a few surplus cells are harder:
with uniform points, 1.4\% on a blocky S-shaped mask (BSP-OT: 9.3\%) and 3.1\% with less than
a row of surplus cells. Starting from a recursive bisection makes results worse.

We also ran the default preset on a billion points of a random walk on one workstation.
It converged in 1\,h\,43\,min using 53\,GiB, and the layout looks right, but at that size no
optimum is known, so this shows that the method scales, not how close it comes.
The exact solver took a median of about two minutes at $10^5$ points. When a proven
optimum is worth that wait, it is the better choice; \ShuffleSnap\ remains the fast answer,
a warm start for the exact solver, and the method for larger sizes and tighter budgets.
The package is on PyPI as \texttt{shufflesnap}; raw results and figure scripts are in a
separate repository.
\label{page:body-end}

\clearpage
\bibliographystyle{plainnat}
\bibliography{refs}

\clearpage
\appendix
\section{Real embeddings}

\panel{fig_visual_mnist}{All 70{,}000 MNIST digits, UMAP layout snapped to a $265\times265$
grid. Colors and graph-paper lines encode input position as in Figure~\ref{fig:progress};
numbers are the mean squared displacement in cells$^2$. Rows: the full grid, the
$64\times64$ box enlarged, each point's displacement, and lattice precision (the fraction of
a point's occupied 8-neighbour cells holding one of its 8 nearest input neighbours).
BSP-OT stays within 3\% of the optimal displacement but halves lattice precision;
RasterFairy breaks the lines at rectangular seams.\label{fig:visual}}
\end{document}
